\section{搜索轨迹与搜索增强建模}

Yuandong Tian 继续展开这条主线，第一步是把搜索过程本身视为监督信号，既保留神经模型的生成能力，又让结果具有可解释的结构。

\subsection{搜索轨迹作为 weak supervision}

讲者展示了搜索增强模型的训练曲线：同样数量级内，搜索轨迹模型能获得 80\% 以上的性能，而纯生成模型在 20\% 附近徘徊。这是因为前者输出的 token 中不仅包含最终计划，还包含 elaborated trace, cost 评估等信息，让 Transformer 有更多上下文。
\begin{figure}[H]
\centering
\includegraphics[width=\textwidth]{frame-search-curve.jpg}
\caption{搜索增强模型与纯生成模型的 training curve，对比 175M 参数的 baseline 与 15M 参数的 trace model。 \protect\footnotemark}
\end{figure}
\footnotetext{视频画面时间区间：00:25:05--00:25:18。}
\begin{knowledgebox}{搜索轨迹的价值}
搜索轨迹让模型在生成每一步之前先 ``思考''：它可以回放执行路径、量化 heuristics、甚至补充静态约束，从而使得 inference 过程中错误率下降、可视化审查变得可行。
\end{knowledgebox}

讲者还强调，在实际实验（Soang、Maze 等）中，训练数据不是空洞的单步指令，而是 solver 生成的 trace、成本值和 plan。通过在 trace 上采样 64 条 plan，再用 heuristics 重排序，模型在少量样本下就能学会长度更长的推理链。

\subsection{Search Former 与 DU Forer}

在此基础上，团队提出 Search Former，并在 DU Forer 中引入 trace drop。通过随机删掉部分搜索 token（如某些 cost 估计或 intermediate node），模型学会在有缺失的情况下依然复原解的核心，类似 self-supervised 的抗干扰训练。
\begin{importantbox}{DU Forer 的训练逻辑}
1）先用 solver 采样完整 trace；2）按级别 drop 特定 token（例如 cost、create、exploration）；3）让 Transformer 根据残余 trace 重建 plan；4）用 loss 反向传播更新 Transformer 和 drop layer。
\end{importantbox}
\begin{enumerate}
    \item \textbf{Level 0：} 保留完整 trace，作为基线输出。
    \item \textbf{Level 1：} 删除 create cost，仅保留 path structure。
    \item \textbf{Level 2：} 删除部分探索动作，强化对关键转折点的建模。
    \item \textbf{Level 3：} 仅保留最终 plan token，模型必须憋出完整 trace。
\end{enumerate}
\begin{figure}[H]
\centering
\includegraphics[width=\textwidth]{frame-dspy-architecture.jpg}
\caption{DU Forer 架构：Search Former 生成 trace，再经过 drop 层、policy head 输出 plan。 \protect\footnotemark}
\end{figure}
\footnotetext{视频画面时间区间：00:32:25--00:32:45。}
\begin{warningbox}{Trace 长度与成本权衡}
长 trace 提供更多上下文，但也需要更高的 Token 长度、GPU 内存以及解析开销；实践中必须控制 drop 策略，避免用数学上完备但推理时不可用的序列。
\end{warningbox}

\subsection{本章小结}

将搜索 trace 作为学习目标，并通过 DU Forer 的 drop 策略调控信息密度，使得小模型也能高效捕捉规划质量。

